% from ega4-8-fr.tex
\oldpage[IV-3]{41}
considérer le cas où $f$ est quasi-projectif. Supposons donc qu'il existe un
$\mathcal O_X$-Module inversible $\mathcal L$ qui soit $f$-ample~; comme $S_0$
est quasi-compact et quasi-séparé il en est de même de $X_0$ (1.2.3), et il y
a donc un $\lambda$ et un $\mathcal O_{X_\lambda}$-Module quasi-cohérent
$\mathcal L_\lambda$ de présentation finie tel que
$\mathcal L=v_\lambda^*(\mathcal L_\lambda)$ (8.5.2, (ii))~; en outre, en
vertu de (8.5.5) on peut supposer $\mathcal L_\lambda$ inversible. Le théorème
dans ce cas est alors conséquence du lemme plus précis suivant.

\begin{lemma}[8.10.5.2]
\label{IV.8.10.5.2-fr}
Supposons $S_0$ quasi-compact, et soit $\mathcal L_\lambda$ un
$\mathcal O_{X_\lambda}$-Module inversible. Pour que $\mathcal L$ soit un
$\mathcal O_X$-Module ample pour $f$ (resp. très ample pour $f$), il faut et
il suffit qu'il existe $\mu\geq\lambda$ tel que $\mathcal L_\mu$ soit ample
pour $f_\mu$ (resp. très ample pour $f_\mu$).
\end{lemma}

La condition étant évidemment suffisante (II, 4.4.10 et 4.6.13), montrons
qu'elle est nécessaire~; les $S_\lambda$ étant quasi-compacts et les
$f_\lambda$ de type fini, on peut, en remplaçant $\mathcal L$ par une puissance
tensorielle convenable, se borner à considérer le cas où $\mathcal L$ est très
ample (II, 4.6.11). En outre, la question étant ici locale sur $S_0$ (vu
(II, 4.4.5) et le fait que $L$ est filtrant), on peut se borner au cas où
$S_0$ (et par suite $S$) est affine. Alors, en vertu de (II, 4.4.1, (ii) et
II, 4.1.2), il existe une $S$-immersion
$j:X\to\mathbf P_S^r=P$ telle que $\mathcal L$ soit isomorphe à
$j^*(\mathcal O_P(1))$. Compte tenu de (8.8.2, (i)), de (ii) et de
(II, 4.1.3), il existe donc un $\mu\geq\lambda$ et une immersion
$j_\mu:X_\mu\to\mathbf P_{S_\mu}^r=P_\mu$ tels que
$j=j_\mu\times1_S$~; utilisant ensuite (II, 4.1.3.2) et (8.5.2.5), on voit
qu'il existe $\nu\geq\mu$ tel que $\mathcal L_\nu$ soit isomorphe à
$j_\nu^*(\mathcal O_{P_\nu}(1))$, ce qui montre que $\mathcal L_\nu$ est très
ample pour $f_\nu$ (II, 4.4.2).
\end{proof}

\subsection{Application aux morphismes quasi-finis}
\label{subsection:IV.8.11-fr}
\label{IV.8.11-fr}

Nous nous proposons dans cette section de démontrer les deux théorèmes
suivants~:

\begin{theorem}[8.11.1]
\EGAReaderTarget{ega4.c04.thm.8-11-1}{IV.8.11.1-fr}
Soit $f:X\to Y$ un morphisme propre, localement de présentation finie, et
quasi-fini. Alors le morphisme $f$ est fini.
\end{theorem}

\begin{theorem}[8.11.2]
\label{IV.8.11.2-fr}
Soit $f:X\to Y$ un morphisme localement de présentation finie, quasi-fini et
séparé. Alors le morphisme $f$ est quasi-affine, et \emph{a fortiori}
quasi-projectif.
\end{theorem}

\begin{remark}[8.11.3]
\label{IV.8.11.3-fr}
On va voir ci-dessous que l'on peut se ramener, pour la démonstration de
(8.11.1) et (8.11.2), au cas où $Y$ est localement noethérien~; on notera que
dans ce cas on obtiendra ainsi une autre démonstration du théorème de
Chevalley (III, 4.4.2).
\end{remark}

\begin{env}[8.11.4]
\label{IV.8.11.4-fr}
Les hypothèses et les conclusions de (8.11.1) et (8.11.2) sont toutes locales
sur $Y$ (I, 6.1, I, 2.6, (II, 5.1.1), (II, 5.4.1) et (II, 6.2.2)), donc on
peut supposer $Y=\operatorname{Spec}(A)$ affine. On sait qu'il existe alors un
sous-anneau $A_0$ de $A$, qui est une $\mathbf Z$-algèbre de type fini, et un
morphisme de type fini $f_0:X_0\to\operatorname{Spec}(A_0)$ tel que $X$
s'identifie à $X_0\otimes_{A_0}A$ et $f$ à $f_0\times1$ (8.9.1). En outre,
$A$ peut être considéré comme limite inductive de ses sous-anneaux contenant
$A_0$ et qui sont des $\mathbf Z$-algèbres de type fini~; utilisant la méthode
de (8.1.2, c)) ainsi que (8.10.5, (v), (xi) et (xii)), on voit qu'il suffit de
démontrer les théorèmes (8.11.1) et (8.11.2) pour $f_0$. Supposons donc
désormais $Y$ noethérien~; utilisant maintenant la méthode de (8.1.2, a))
ainsi que (8.10.5, (v), (ix), (x), (xi) et (xii)) on peut remplacer $Y$ par
$\operatorname{Spec}(\mathcal O_y)$, où $y$ est un
